\documentclass[12pt]{extarticle}
\usepackage[utf8]{inputenc}
% \usepackage{cite}
\usepackage{geometry}
\usepackage{hyperref}
\hypersetup{
    colorlinks=true,
    citecolor=blue,
    linkcolor=black,
    urlcolor=blue
}
\geometry{margin=1in}

\title{Small Models, Big Reactions: Studying Efficient Emoji Prediction}
\author{Minh Vu, Tri Than, Duy Nguyen}
\date{\today}

\begin{document}
\maketitle

% \begin{abstract}
% Provide a brief summary of the proposal, including the core objective, methodology, and expected impact.
% \end{abstract}

\section{Problem and Motivation}

Emoji are widely used in online communication to express emotion, tone, humor, 
and reactions that may not be obvious from text alone. We study emoji prediction 
as a multi-class text classification task: given a short social-media post with 
its emoji removed, the model predicts one of a fixed set of emoji classes.

Prior work has shown that both simple classifiers and pretrained language models
can perform this task, but these approaches vary greatly in predictive performance
and computational cost. Understanding this tradeoff is useful because a smaller
model may be preferable when computational resources or inference speed are limited,
even if it sacrifices some predictive performance. We therefore ask:
\textbf{What is the tradeoff between model complexity and emoji-prediction performance?}
Using the TweetEval emoji benchmark, we will compare models of increasing complexity,
including TF-IDF + logistic regression, a BiLSTM classifier, DistilBERT, and a larger
Transformer such as BERT-base or RoBERTa-base. We will evaluate macro-F1, parameter
count, and inference latency, and analyze frequently confused emoji classes. This
study will help us understand whether the performance gains of larger models justify
their additional computational cost, and whether those gains are concentrated in
difficult or semantically similar emoji classes.

\section{Literature Survey}

Emoji prediction has been studied as a social-media text classification task in prior 
work. The SemEval-2018 Task 2 benchmark formalized the problem as predicting the most 
likely emoji associated with a tweet and evaluated systems using macro-F1 because of 
the imbalance among emoji classes \cite{barbieri2018semeval}. A wide range of approaches 
were explored in the shared task, including bag-of-words classifiers, support vector 
machines, convolutional neural networks, recurrent neural networks, and attention-based 
models. Some relatively simple systems were surprisingly competitive, while more complex 
neural models attempted to capture semantic and contextual information more effectively.

Later, TweetEval incorporated emoji prediction into a standardized benchmark for Twitter
-specific classification tasks and provided fixed training, validation, and test splits 
together with stronger pretrained language-model baselines \cite{barbieri2020tweeteval}. 
TweetEval showed that pretrained Transformer models can improve performance, but emoji 
prediction remains difficult because the label distribution is highly skewed and several 
emoji classes have similar meanings. These results motivate our study of whether the 
additional complexity of larger models produces enough improvement to justify their higher 
computational cost.

\section{Data Requirements}

We will use the TweetEval emoji classification benchmark
\cite{barbieri2020tweeteval}, which contains short social-media posts labeled with one
of 20 emojis. The provided splits contain approximately 45,000 training, 5,000
validation, and 50,000 test examples. The dataset directly matches our task and enables
comparison with established baselines.

TweetEval contains substantial class imbalance and semantically similar emoji classes,
making macro-F1 and class-level error analysis important. We will preserve the provided
splits and apply only limited preprocessing, such as removing residual Unicode artifacts
when necessary.

\section{Experiment Plan}

We will compare several models of increasing complexity on the TweetEval emoji prediction task. 
Our baseline will be TF-IDF features with logistic regression. We will then compare this baseline 
with a BiLSTM classifier, DistilBERT, and a larger pretrained Transformer such as BERT-base or 
RoBERTa-base. All models will use the same train, validation, and test splits to ensure a fair 
comparison.

We will evaluate predictive performance using macro-F1, which is appropriate for the imbalanced 
emoji classes. To measure computational cost, we will record parameter count and inference latency 
under the same hardware and batch settings. We will also perform class-level and qualitative error 
analysis using confusion matrices and representative misclassified examples, focusing on whether 
larger models better distinguish semantically similar emoji classes.

\section{Feasibility and Next Steps}

We will run experiments in Google Colab Pro, using an NVIDIA T4 GPU as our expected
resource and an A100 when available. Our first experiment will establish the TF-IDF +
logistic regression baseline and verify the TweetEval training and evaluation pipeline.
By the milestone, we aim to evaluate the baseline, BiLSTM, and at least one pretrained
Transformer, with preliminary macro-F1, model-size, and inference-latency results. By
the final report, we will complete the model comparison and qualitative error analysis.

Potential challenges include GPU availability, class imbalance, noisy social-media text,
and preprocessing artifacts. If compute becomes limiting, we will reduce the number of
larger Transformer models while preserving the performance--complexity comparison.

\bibliographystyle{plain}
\bibliography{references}

\end{document}
